\begin{examplebox}

	\textbf{\textcolor{CExample}{例 1（基础）}}

	\textbf{\textcolor{CExample}{题目：}}
	求 $S_4 = 1\cdot1!+2\cdot2!+3\cdot3!+4\cdot4!$ 的值。

	\textbf{\textcolor{CExample}{解题步骤：}}
	\begin{description}[style=nextline, leftmargin=2.8em, labelsep=0.5em, itemsep=0.45em]
		\item[\textbf{第一步（裂项转化）：}] 对每一项应用公式 $k\cdot k! = (k+1)!-k!$。
			\par\textbf{理由：} 由公式 $n\cdot n! = (n+1)!-n!$ 直接得到。
		\item[\textbf{第二步（展开和式）：}] 写出 $S_4 = (2!-1!)+(3!-2!)+(4!-3!)+(5!-4!)$。
			\par\textbf{理由：} 将 $k=1,2,3,4$ 分别代入裂项公式。
		\item[\textbf{第三步（消项求值）：}] 合并得
			\[
				\begin{aligned}
					S_4 &= 5! - 1! \\
					&= 120 - 1 \\
					&= 119.
			\end{aligned}
			\]
			\par\textbf{理由：} 相邻项 $-1!$ 与 $+2!$ 的 $2!$ 抵消，依次类推，只剩 $5!$ 和 $-1!$，且 $1! = 1$，$5! = 120$。
	\end{description}

	\textbf{\textcolor{CExample}{关键结论：}}
	$S_4 = 119$，验证了公式 $S_n = (n+1)!-1$。

	\medskip
	\textbf{\textcolor{CExample}{例 2（稍变形）}}

	\textbf{\textcolor{CExample}{题目：}}
	求 $T_n = \sum_{k=2}^{n} k\cdot k!$ 的表达式。

	\textbf{\textcolor{CExample}{解题步骤：}}
	\begin{description}[style=nextline, leftmargin=2.8em, labelsep=0.5em, itemsep=0.45em]
		\item[\textbf{第一步（裂项转化）：}] 每项 $k\cdot k! = (k+1)!-k!$。
			\par\textbf{理由：} 同例1。
		\item[\textbf{第二步（展开和式）：}]
			\[
				\begin{aligned}
					T_n &= \sum_{k=2}^{n} [(k+1)!-k!] \\
					&= (3!-2!)+(4!-3!)+\cdots+((n+1)!-n!).
			\end{aligned}
			\]
			\par\textbf{理由：} 代入裂项并展开。
		\item[\textbf{第三步（消项化简）：}] 消去中间项得
			\[
				\begin{aligned}
					T_n &= (n+1)! - 2! \\
					&= (n+1)! - 2.
			\end{aligned}
			\]
			\par\textbf{理由：} 首项剩余 $-2!$，末项保留 $(n+1)!$，其余抵消。
	\end{description}

	\textbf{\textcolor{CExample}{关键结论：}}
	$T_n = (n+1)! - 2$。

	\medskip
	\textbf{\textcolor{CExample}{例 3（综合应用）}}

	\textbf{\textcolor{CExample}{题目：}}
	用数学归纳法证明 $1\cdot1!+2\cdot2!+\cdots+n\cdot n! = (n+1)! - 1$ 对所有正整数 $n$ 成立。

	\textbf{\textcolor{CExample}{解题步骤：}}
	\begin{description}[style=nextline, leftmargin=2.8em, labelsep=0.5em, itemsep=0.45em]
		\item[\textbf{第一步（基础验证）：}] 当 $n=1$ 时，左式 $1\cdot1! = 1$，右式 $2!-1 = 1$，等式成立。
			\par\textbf{理由：} 直接计算 $n=1$ 情况。
		\item[\textbf{第二步（归纳假设）：}] 假设 $n=k$ 时成立，即 $\sum_{i=1}^{k} i\cdot i! = (k+1)! - 1$。
			\par\textbf{理由：} 归纳法假设。
		\item[\textbf{第三步（归纳递推）：}] 当 $n=k+1$ 时，
			\[
				\begin{aligned}
					\text{左} &= \sum_{i=1}^{k} i\cdot i! + (k+1)\cdot(k+1)! \\
					&= \bigl[(k+1)! - 1\bigr] + (k+1)\cdot(k+1)! \quad\text{(代入归纳假设)}\\
					&= (k+1)! \cdot [1 + (k+1)] - 1 \\
					&= (k+1)! \cdot (k+2) - 1 \\
					&= (k+2)! - 1.
			\end{aligned}
			\]
			即得到右式，故等式成立。
			\par\textbf{理由：} 利用假设加上新项，提取公因式化简，本质是裂项公式 $ (k+1)\cdot(k+1)! = (k+2)! - (k+1)!$ 的应用。
	\end{description}

	\textbf{\textcolor{CExample}{关键结论：}}
	等式得证，同时验证了裂项求和公式的正确性。

\end{examplebox}
